\documentclass[11pt,a4paper]{article}
\usepackage{amsmath,amssymb,geometry,booktabs,hyperref,pifont}
\geometry{margin=2.5cm}

\newcommand{\Nf}{N_{\!f}}
\newcommand{\Nc}{N_{\!c}}
\newcommand{\rep}[1]{\mathbf{#1}}
\newcommand{\conjrep}[1]{\overline{\mathbf{#1}}}
\newcommand{\Tr}{\mathrm{Tr}}

\title{Meson decomposition from the ISS rank condition:\\
four starting points and SM charge assignment}
\author{A.~Rivero\\[2pt]
\small\textit{BIFI, Universidad de Zaragoza, E-50018 Zaragoza, Spain}\\
\small\texttt{arivero@unizar.es}}
\date{}

\begin{document}
\maketitle

\begin{abstract}
In the Intriligator--Seiberg--Shih (ISS) model of metastable SUSY
breaking, the rank condition of the meson matrix $M = \tilde Q Q$
splits $\Nf$ flavors into $k_u = \Nf - \Nc$ massless pseudo-moduli
and $k_d = \Nc$ tree-level massive directions.
We tabulate the meson decomposition under
$\mathrm{SU}(k_u)\times\mathrm{SU}(k_d)\times\mathrm{U}(1)$
for four choices of $(\Nc,\Nf)$ within or near the conformal window:
$(3,6)$, $(3,5)$, $(2,6)$, and $(2,5)$.
For each case we examine how Standard Model electric charges can be
assigned to the constituent quarks and how the resulting meson
quantum numbers compare with the SM spectrum.
The case $(\Nc,\Nf)=(3,5)$, which is the sBootstrap starting point,
uniquely matches the SM charge assignment with two up-type and three
down-type quarks.
The alternative $(2,5)$ reproduces the SM non-abelian gauge group
$\mathrm{SU}(3)\times\mathrm{SU}(2)$ from the flavor decomposition,
but with an inverted mass--charge correspondence.
We discuss the tension between these two viewpoints and the role
of the $\Nf=6\to 5$ transition (top decoupling) in connecting them.
\end{abstract}

%% ===================================================================
\section{Introduction}
\label{sec:intro}
%% ===================================================================

The ISS mechanism~\cite{ISS2006} provides a calculable model of
metastable supersymmetry breaking in the magnetic dual of SQCD.
A central ingredient is the \emph{rank condition}: the magnetic
meson vacuum expectation value has rank at most $\Nf - \Nc$, so that
$\Nc$ of the $\Nf$ flavor directions acquire tree-level masses while
$\Nf - \Nc$ directions remain as massless pseudo-moduli.
This split is a robust consequence of the duality structure and
does not depend on details of the K\"ahler potential.

From a phenomenological perspective, the $\Nf\times\Nf$ meson matrix
$M_{ij} = \tilde Q_i Q_j$ of Seiberg duality~\cite{Seiberg1995}
transforms as
$(\Nf^2 - 1) \oplus 1$ under $\mathrm{SU}(\Nf)$, with $\Nf^2$
components in total.
The rank condition selects a vacuum that breaks
$\mathrm{SU}(\Nf)$ to
$\mathrm{SU}(k_u)\times\mathrm{SU}(k_d)\times\mathrm{U}(1)$,
where $k_u = \Nf - \Nc$ and $k_d = \Nc$.
The decomposition of the $\Nf^2$ mesons under this residual symmetry
depends on $(\Nc, \Nf)$ and determines the low-energy spectrum.

Different values of $(\Nc, \Nf)$ within the Seiberg conformal
window $\Nc + 1 \le \Nf < 3\Nc$ lead to different meson
decompositions and different prospects for embedding the Standard
Model.  In this note we systematically compare four choices:
$(\Nc,\Nf) = (3,6)$, $(3,5)$, $(2,6)$, and $(2,5)$.
We then address how Standard Model electric charges can be
distributed among the $\Nf$ quarks in each case, and examine which
starting point most naturally reproduces the SM meson and Higgs
content.

Throughout, our conventions for Seiberg duality follow~\cite{Seiberg1995},
eqs.~(1.1)--(1.2) for the electric theory and eqs.~(3.1)--(3.3) for the
magnetic dual.  The ISS superpotential and vacuum structure
follow~\cite{ISS2006}, eqs.~(2.4)--(2.6).

%% ===================================================================
\section{Setup: Seiberg duality and the ISS mechanism}
\label{sec:setup}
%% ===================================================================

\subsection{Seiberg duality}

Consider $\mathcal{N}=1$ SQCD with gauge group $\mathrm{SU}(\Nc)$
and $\Nf$ flavors of quarks $Q^i$, $\tilde Q_j$ in the fundamental
and antifundamental representations
($i,j = 1,\ldots,\Nf$)~\cite{Seiberg1995}.
The global symmetry is
\begin{equation}
  \mathrm{SU}(\Nf)_L \times \mathrm{SU}(\Nf)_R
  \times \mathrm{U}(1)_B \times \mathrm{U}(1)_R\,.
  \label{eq:global-sym}
\end{equation}
The gauge-invariant composite operators include the meson
\begin{equation}
  M^i{}_j = \tilde Q^i \, Q_j\,,
  \label{eq:meson-def}
\end{equation}
which transforms as $(\conjrep{\Nf}, \rep{\Nf})$ under
$\mathrm{SU}(\Nf)_L \times \mathrm{SU}(\Nf)_R$, i.e.\ as a
bifundamental with $\Nf^2$ components.

In the conformal window, $\Nc + 1 \le \Nf < 3\Nc$, the theory
admits a Seiberg dual description\footnote{The ISS one-loop
analysis is perturbatively controlled in the narrower
\emph{free magnetic range} $\Nc + 1 \le \Nf \le \frac{3}{2}\Nc$.
All four cases studied here have $\Nf > \frac{3}{2}\Nc$ and
therefore sit in the conformal window but outside the free
magnetic range.  The rank condition itself is a tree-level
statement valid throughout the conformal window; only the
Coleman--Weinberg stabilization of pseudo-moduli requires
the free magnetic range for perturbative control.}: an $\mathrm{SU}(N)$ gauge theory
with $N = \Nf - \Nc$, containing $\Nf$ dual quarks $\varphi_i$,
$\tilde\varphi^j$ and an elementary gauge-singlet field $\Phi^i{}_j$
identified with the meson $M^i{}_j$.
The dual superpotential is~\cite{Seiberg1995}
\begin{equation}
  W_{\text{dual}} = h\,\Tr(\varphi\,\Phi\,\tilde\varphi)\,,
  \label{eq:dual-W}
\end{equation}
where $h$ is the dual Yukawa coupling (cf.\ eq.~(3.3) of~\cite{Seiberg1995}).

\subsection{ISS metastable SUSY breaking}

Adding a mass term $W_{\text{mass}} = \Tr(m\, M)$ in the electric
theory maps in the dual to a linear term for $\Phi$.
The full magnetic superpotential becomes~\cite{ISS2006}
\begin{equation}
  W = h\,\Tr(\varphi\,\Phi\,\tilde\varphi)
    - h\mu^2 \Tr\Phi\,,
  \label{eq:ISS-W}
\end{equation}
where $\mu^2$ encodes the electric quark mass
(eq.~(2.4) of~\cite{ISS2006}).
The $F$-term conditions $F_\Phi = 0$ require
\begin{equation}
  \varphi\,\tilde\varphi = \mu^2\,\mathbf{1}_{\Nf}\,,
  \label{eq:F-flat}
\end{equation}
but $\varphi$ and $\tilde\varphi$ are $N\times\Nf$ matrices, so
$\varphi\,\tilde\varphi$ has rank at most $N = \Nf - \Nc$.
Therefore eq.~\eqref{eq:F-flat} cannot be satisfied for all $\Nf$
flavors simultaneously: this is the \emph{rank condition}
(following from the $F$-term conditions of the
superpotential~(1.3) of~\cite{ISS2006}).

The metastable vacuum takes the form (eq.~(2.6) of~\cite{ISS2006})
\begin{equation}
  \varphi_0 = \tilde\varphi_0^T
  = \begin{pmatrix} \mu\,\mathbf{1}_N \\ 0_{(\Nf-N)\times N} \end{pmatrix},
  \qquad
  \Phi_0 \;\text{arbitrary (pseudo-moduli)}\,.
  \label{eq:ISS-vacuum}
\end{equation}
This splits the $\Nf$ flavors into two groups:
\begin{itemize}
\item \textbf{$k_u = N = \Nf - \Nc$ directions:} the $F$-terms are
  satisfied, $\varphi\tilde\varphi = \mu^2\mathbf{1}_N$.
  At tree level the meson components $\Phi$ along these directions
  are pseudo-moduli (massless).  They acquire a mass only at one loop
  through the Coleman--Weinberg potential.
\item \textbf{$k_d = \Nc$ directions:} the $F$-terms are not
  satisfied, $F_\Phi \ne 0$.  The meson components along these
  directions have tree-level masses of order $h\mu$.
\end{itemize}
The vacuum energy is $V_{\min} = (k_d)\,|h^2\mu^4|$
(eq.~(2.5) of~\cite{ISS2006}).

\subsection{Meson decomposition}

The vacuum~\eqref{eq:ISS-vacuum} breaks the $\mathrm{SU}(\Nf)$
flavor symmetry to
\begin{equation}
  \mathrm{SU}(\Nf) \;\longrightarrow\;
  \mathrm{SU}(k_u) \times \mathrm{SU}(k_d)
  \times \mathrm{U}(1)\,.
  \label{eq:breaking-pattern}
\end{equation}
The $\Nf^2$ meson components decompose as
\begin{equation}
  \Nf^2 \;=\;
  \underbrace{(k_u^2-1)+1}_{\mathrm{SU}(k_u)\text{ sector}}
  \;+\;
  \underbrace{(k_d^2-1)+1}_{\mathrm{SU}(k_d)\text{ sector}}
  \;+\;
  2\,k_u\, k_d\,,
  \label{eq:meson-counting}
\end{equation}
where the last term accounts for the bifundamental and its conjugate.
In representation language:
\begin{equation}
  \Nf^2 \;\to\;
  (\rep{adj},\rep{1})_0
  \oplus (\rep{1},\rep{adj})_0
  \oplus (\rep{1},\rep{1})_0
  \oplus (\rep{1},\rep{1})_0
  \oplus (\rep{k_u},\conjrep{k_d})_{+1}
  \oplus (\conjrep{k_u},\rep{k_d})_{-1}\,,
  \label{eq:branching}
\end{equation}
where the two singlets arise from the two independent
$\mathrm{U}(1)$ projections of the trace.

The physical content depends on $k_u$ and $k_d$, i.e.\ on the
choice of $(\Nc,\Nf)$.


%% ===================================================================
\section{The four decompositions}
\label{sec:four}
%% ===================================================================

We now tabulate the four cases.  Table~\ref{tab:summary} provides a
compact comparison; the subsections below give the details.

\begin{table*}[t]
\centering
\caption{Meson decomposition for four $(\Nc,\Nf)$ starting points.
  $N = \Nf - \Nc$ is the dual gauge group rank.
  ``CW'' indicates conformal window status.
  The branching is under
  $\mathrm{SU}(k_u)\times\mathrm{SU}(k_d)\times\mathrm{U}(1)$.}
\label{tab:summary}
\smallskip
\begin{tabular}{@{}cccccccl@{}}
\toprule
$\Nc$ & $\Nf$ & $N$ & $k_u$ & $k_d$ & Mesons & CW &
Branching rule \\
\midrule
3 & 6 & 3 & 3 & 3 & 36 & $4\!\le\!6\!<\!9$ \checkmark &
$(\rep{8},\rep{1})_0 \oplus (\rep{1},\rep{8})_0
 \oplus 2\times(\rep{1},\rep{1})_0
 \oplus (\rep{3},\conjrep{3})_{+1}
 \oplus (\conjrep{3},\rep{3})_{-1}$  \\[4pt]
3 & 5 & 2 & 2 & 3 & 25 & $4\!\le\!5\!<\!9$ \checkmark &
$(\rep{3},\rep{1})_0 \oplus (\rep{1},\rep{8})_0
 \oplus 2\times(\rep{1},\rep{1})_0
 \oplus (\rep{2},\conjrep{3})_{+1}
 \oplus (\conjrep{2},\rep{3})_{-1}$  \\[4pt]
2 & 6 & 4 & 4 & 2 & 36 & $3\!\le\!6\!\not<\!6$ \ding{55}\!\! &
$(\rep{15},\rep{1})_0 \oplus (\rep{1},\rep{3})_0
 \oplus 2\times(\rep{1},\rep{1})_0
 \oplus (\rep{4},\conjrep{2})_{+1}
 \oplus (\conjrep{4},\rep{2})_{-1}$  \\[4pt]
2 & 5 & 3 & 3 & 2 & 25 & $3\!\le\!5\!<\!6$ \checkmark &
$(\rep{8},\rep{1})_0 \oplus (\rep{1},\rep{3})_0
 \oplus 2\times(\rep{1},\rep{1})_0
 \oplus (\rep{3},\conjrep{2})_{+1}
 \oplus (\conjrep{3},\rep{2})_{-1}$  \\
\bottomrule
\end{tabular}
\end{table*}

\subsection{$(\Nc,\Nf) = (3,6)$: symmetric split}
\label{sec:36}

The dual gauge group is $\mathrm{SU}(3)$, and the theory is
self-dual in the sense that both the electric and magnetic theories
have the same gauge group rank.
The conformal window condition is $\Nc + 1 = 4 \le 6 < 9 = 3\Nc$,
which is satisfied.

With $k_u = k_d = 3$ the flavor symmetry breaks as
\begin{equation}
  \mathrm{SU}(6) \;\to\;
  \mathrm{SU}(3)_u \times \mathrm{SU}(3)_d \times \mathrm{U}(1)\,.
\end{equation}
The adjoint of $\mathrm{SU}(6)$ decomposes as
\begin{equation}
  \rep{35} \;\to\;
  (\rep{8},\rep{1})_0 \oplus (\rep{1},\rep{8})_0
  \oplus (\rep{1},\rep{1})_0
  \oplus (\rep{3},\conjrep{3})_{+1}
  \oplus (\conjrep{3},\rep{3})_{-1}\,.
  \label{eq:adj-36}
\end{equation}
Adding the singlet from the full $\Nf^2 = 36$ gives
$8 + 8 + 1 + 1 + 9 + 9 = 36$.

This is the starting point of the Cacciapaglia--Sannino dualised
Standard Model~\cite{CacciapagliaSannino2024}, which embeds six
quark flavors into a Seiberg-dual framework and obtains an
18-Higgs-doublet structure.
The symmetric split $k_u = k_d = 3$ means that no preference is
built in for which flavors are ``up-type'' and which are
``down-type.''

\subsubsection*{Top decoupling: from $(3,6)$ to $(3,5)$}

The $k_u = k_d$ symmetry is broken by giving one flavor (the top)
a large mass $m_t$.  In the electric theory, the heavy top
is integrated out below the scale $m_t$, reducing $\Nf$ from 6
to~5.  Re-deriving the Seiberg dual for the resulting
$\mathrm{SU}(3)$ theory with $\Nf = 5$ gives a new magnetic dual
with gauge group $\mathrm{SU}(\Nf - \Nc) = \mathrm{SU}(2)$:
\begin{equation}
  \mathrm{SU}(3)_{\text{dual}}^{(\Nf=6)}
  \;\xrightarrow{\;\text{top decoupling}\;}
  \mathrm{SU}(2)_{\text{dual}}^{(\Nf=5)}\,.
  \label{eq:decoupling-36}
\end{equation}
Below the scale $m_t$, the effective theory has $\Nf = 5$ active
flavors with dual gauge group $\mathrm{SU}(2)$.  The $6\times 6$
meson matrix reduces to $5\times 5 = 25$ components, and the ISS
rank condition now gives the asymmetric split $k_u = 2$,
$k_d = 3$: this is the $(3,5)$ case of
Section~\ref{sec:35}.

The 36 mesons of the parent $(3,6)$ theory decompose into
25~mesons of the $(3,5)$ daughter theory plus 11~states involving
the top flavor (5~charged under the residual $\mathrm{SU}(5)$ plus
6~top-mixed states), which become heavy and decouple.
The asymmetric $k_u = 2$, $k_d = 3$ split is thus a
consequence of the large top mass (which is an input, not a
prediction), rather than being imposed by hand on the flavor
structure.

\subsection{$(\Nc,\Nf) = (3,5)$: the sBootstrap case}
\label{sec:35}

The dual gauge group is $\mathrm{SU}(2)$.
The conformal window condition $4 \le 5 < 9$ is satisfied.

With $k_u = 2$ and $k_d = 3$ the flavor symmetry breaks as
\begin{equation}
  \mathrm{SU}(5) \;\to\;
  \mathrm{SU}(2)_u \times \mathrm{SU}(3)_d \times \mathrm{U}(1)\,.
\end{equation}
The adjoint decomposes as
\begin{equation}
  \rep{24} \;\to\;
  (\rep{3},\rep{1})_0 \oplus (\rep{1},\rep{8})_0
  \oplus (\rep{1},\rep{1})_0
  \oplus (\rep{2},\conjrep{3})_{+1}
  \oplus (\conjrep{2},\rep{3})_{-1}\,.
  \label{eq:adj-35}
\end{equation}
Adding the singlet: $3 + 8 + 1 + 1 + 6 + 6 = 25$.

This is the case relevant to the sBootstrap
program~\cite{Rivero2024}: the $k_u = 2$ massless directions are
identified with the two up-type quarks $(u,c)$, and the $k_d = 3$
massive directions with the three down-type quarks $(d,s,b)$.
The 25 mesons organize as open-flavor states in
$(\rep{2},\conjrep{3})$ plus their conjugates and flavor-diagonal
states in $(\rep{3},\rep{1})$ and $(\rep{1},\rep{8})$.

\subsection{$(\Nc,\Nf) = (2,6)$: Pati--Salam-like}
\label{sec:26}

The dual gauge group would be $\mathrm{SU}(4)$.
However, the conformal window condition
$\Nc + 1 = 3 \le 6 < 6 = 3\Nc$ \emph{fails}: $\Nf = 3\Nc$
is the boundary of asymptotic freedom, and for $\Nf \ge 3\Nc$ the
theory is infrared-free.
This case therefore sits outside (or at the edge of) the regime
where Seiberg duality and the ISS mechanism are under perturbative
control.

With $k_u = 4$ and $k_d = 2$ the formal decomposition is
\begin{equation}
  \mathrm{SU}(6) \;\to\;
  \mathrm{SU}(4)_u \times \mathrm{SU}(2)_d \times \mathrm{U}(1)\,.
\end{equation}
The adjoint decomposes as
\begin{equation}
  \rep{35} \;\to\;
  (\rep{15},\rep{1})_0 \oplus (\rep{1},\rep{3})_0
  \oplus (\rep{1},\rep{1})_0
  \oplus (\rep{4},\conjrep{2})_{+1}
  \oplus (\conjrep{4},\rep{2})_{-1}\,.
  \label{eq:adj-26}
\end{equation}
Adding the singlet: $15 + 3 + 1 + 1 + 8 + 8 = 36$.

The $\mathrm{SU}(4) \times \mathrm{SU}(2)$ structure is
reminiscent of the Pati--Salam group.
If $\mathrm{SU}(4)$ is further broken to
$\mathrm{SU}(3) \times \mathrm{U}(1)$,
the $\mathrm{U}(1)$ factor plays the role of a
lepton-number symmetry.
However, the failure of the conformal window condition makes this
case less attractive as a controlled starting point.

\subsection{$(\Nc,\Nf) = (2,5)$: SM gauge structure}
\label{sec:25}

The dual gauge group is $\mathrm{SU}(3)$.
The conformal window condition $3 \le 5 < 6$ is satisfied.

With $k_u = 3$ and $k_d = 2$ the flavor symmetry breaks as
\begin{equation}
  \mathrm{SU}(5) \;\to\;
  \mathrm{SU}(3)_u \times \mathrm{SU}(2)_d \times \mathrm{U}(1)\,.
  \label{eq:SM-group}
\end{equation}
The adjoint decomposes as
\begin{equation}
  \rep{24} \;\to\;
  (\rep{8},\rep{1})_0 \oplus (\rep{1},\rep{3})_0
  \oplus (\rep{1},\rep{1})_0
  \oplus (\rep{3},\conjrep{2})_{+1}
  \oplus (\conjrep{3},\rep{2})_{-1}\,.
  \label{eq:adj-25}
\end{equation}
Adding the singlet: $8 + 3 + 1 + 1 + 6 + 6 = 25$.

The residual symmetry~\eqref{eq:SM-group} is precisely the SM
non-abelian gauge group $\mathrm{SU}(3)_c \times \mathrm{SU}(2)_L$.
The bifundamental $(\rep{3},\conjrep{2})$ is exactly the
representation of a left-handed quark doublet.
This match between a flavor decomposition and the SM gauge
structure arises within the general Seiberg duality
framework~\cite{Seiberg1995} and has been discussed in the
context of the dualised SM~\cite{CacciapagliaSannino2024}.


%% ===================================================================
\section{Standard Model charge assignment}
\label{sec:charges}
%% ===================================================================

In the Standard Model, the six quark flavors carry electric charges
$Q_{\text{em}} = +2/3$ (up-type: $u, c, t$) and
$Q_{\text{em}} = -1/3$ (down-type: $d, s, b$).
With only five light flavors (excluding the top), we have 2 up-type
and 3 down-type quarks.

The ISS rank condition assigns the $k_u$ flavors to the massless
(pseudo-moduli) sector and the $k_d$ flavors to the massive sector.
We now examine whether the SM charge assignment is compatible with
this split for each $(\Nc,\Nf)$.

\subsection{Charge assignment for $(\Nc,\Nf) = (3,5)$}

Here $k_u = 2$ and $k_d = 3$.
Identifying the 2 massless directions with the up-type quarks
$(u,c)$ and the 3 massive directions with the down-type quarks
$(d,s,b)$ gives a direct match with the SM:
\begin{center}
\begin{tabular}{cccl}
\toprule
Sector & $k$ & $Q_{\text{em}}$ & Flavors \\
\midrule
Massless (pseudo-moduli) & 2 & $+2/3$ & $u,\; c$ \\
Massive (tree-level) & 3 & $-1/3$ & $d,\; s,\; b$ \\
\bottomrule
\end{tabular}
\end{center}
The mesons in the bifundamental
$(\rep{2},\conjrep{3})$ carry
$Q_{\text{em}} = (+2/3) - (-1/3) = +1$ and
$(\conjrep{2},\rep{3})$ carry $Q_{\text{em}} = -1$,
while the diagonal sectors $(\rep{3},\rep{1})$ and
$(\rep{1},\rep{8})$ are electrically neutral.
The full charge spectrum of the 25 mesons is therefore
\begin{equation}
  12_{\,0} + 6_{+1} + 6_{-1} + 1_{\,0} = 25\,,
  \label{eq:meson-charges-35}
\end{equation}
where the subscript denotes $Q_{\text{em}}$.
Each meson $M_{ij} = \bar q_i\, q_j$ carries charge
$Q_{\text{em}}(M_{ij}) = -Q_i + Q_j$, with $Q_i$ the electric
charge of quark flavor~$i$.
The thirteen neutral mesons ($12 + 1$) correspond to
$(\rep{3},\rep{1})_0 \oplus (\rep{1},\rep{8})_0
\oplus 2\times(\rep{1},\rep{1})_0$
and the twelve charged ones to
$(\rep{2},\conjrep{3})_{+1} \oplus (\conjrep{2},\rep{3})_{-1}$.

This charge assignment works because the ISS split $k_u = 2$,
$k_d = 3$ exactly matches the SM flavor counting.

\subsection{Charge assignment for $(\Nc,\Nf) = (2,5)$}

Here $k_u = 3$ and $k_d = 2$.
The ISS rank condition assigns 3 flavors to the massless sector and
2 to the massive sector.  If we want to identify these with SM
quarks, we face an inversion: 3 up-type vs.\ 2 down-type, which is
the opposite of the SM counting.

One could instead assign the 3 massless directions to the down-type
quarks $(d,s,b)$ and the 2 massive ones to the up-type quarks
$(u,c)$, but then the ISS mass hierarchy is inverted relative to the
SM: in the SM, the up-type quarks are heavier (on average) than the
down-type quarks in each generation.  More fundamentally, the massive
ISS directions are the ones whose $F$-terms are not satisfied ---
they ``feel'' the SUSY breaking directly.

\subsection{Charge assignment for $(\Nc,\Nf) = (3,6)$}

Here $k_u = k_d = 3$, so the split is symmetric.
The SM charge assignment requires distinguishing $u$- and $d$-type
quarks, but the ISS vacuum treats both sectors identically at tree
level.  To reach the SM content one must invoke additional dynamics
--- for instance, decoupling the top quark to reduce
$\Nf = 6 \to 5$ --- which breaks the $k_u = k_d$ symmetry.

This is the starting point of the Cacciapaglia--Sannino
construction~\cite{CacciapagliaSannino2024}, in which the full
six-flavor theory is dualised and the SM emerges after appropriate
symmetry-breaking steps.

\subsection{Charge assignment for $(\Nc,\Nf) = (2,6)$}

Here $k_u = 4$ and $k_d = 2$, and the theory sits at the edge of
the conformal window.  Neither $k_u = 4$ nor $k_d = 2$ matches
the SM flavor counting directly (we need 2 up-type and 3 down-type,
or at most 3 and 3 for the full six-flavor case).  This case does
not admit a natural SM charge assignment without further breaking of
$\mathrm{SU}(4)$.

Table~\ref{tab:charges} summarizes the charge-assignment analysis.

\begin{table}[t]
\centering
\caption{SM charge assignment compatibility for each
  $(\Nc,\Nf)$.}
\label{tab:charges}
\smallskip
\begin{tabular}{@{}cccccl@{}}
\toprule
$\Nc$ & $\Nf$ & $k_u$ & $k_d$ & Match? & Comment \\
\midrule
3 & 6 & 3 & 3 & Partial &
  Symmetric; needs top decoupling \\
3 & 5 & 2 & 3 & \textbf{Yes} &
  $k_u = n_u$, $k_d = n_d$ exactly \\
2 & 6 & 4 & 2 & No &
  Outside conformal window \\
2 & 5 & 3 & 2 & Inverted &
  $k_u \leftrightarrow k_d$ vs.\ SM \\
\bottomrule
\end{tabular}
\end{table}


%% ===================================================================
\section{Electroweak embedding}
\label{sec:ew}
%% ===================================================================

The comparison between the $(3,5)$ and $(2,5)$ cases reveals a
tension between two desiderata: \emph{charge assignment} and
\emph{gauge structure}.

\subsection{The $(3,5)$ perspective: charge assignment}

For $(\Nc,\Nf) = (3,5)$, the ISS split produces
$\mathrm{SU}(2)_u \times \mathrm{SU}(3)_d \times \mathrm{U}(1)$.
The $\mathrm{SU}(2)$ factor acts on the $k_u = 2$ massless
directions, i.e.\ the up-type quarks.
The mesons in $(\rep{2},\conjrep{3})_{+1}$ are electroweak doublets
of color anti-triplets --- exactly the quantum numbers of a Higgs
doublet in the dualised SM.

The identification of this $\mathrm{SU}(2)$ with
$\mathrm{SU}(2)_L$ is natural from the ISS perspective: the
massless pseudo-moduli are the directions along which the meson VEV
can fluctuate freely, and in the SM the left-handed doublets are
massless before electroweak symmetry breaking.
The $\mathrm{SU}(2)$ here rotates the massless modes, paralleling
$\mathrm{SU}(2)_L$ rotating the left-handed doublets.

This is the sBootstrap viewpoint~\cite{Rivero2024}: the mesons
$M_{ij} = \tilde q_i q_j$ with one up-type and one down-type index
carry the quantum numbers of the scalar partners of SM quarks
(squarks/sfermions), and can be identified with pseudoscalar mesons in
the hadronic spectrum.

\subsection{The $(2,5)$ perspective: gauge structure}

For $(\Nc,\Nf) = (2,5)$, the ISS split produces
$\mathrm{SU}(3)_u \times \mathrm{SU}(2)_d \times \mathrm{U}(1)$
--- exactly the SM non-abelian gauge group.
The $(\rep{3},\conjrep{2})$ bifundamental is exactly a left-handed
quark doublet representation.
The $(\rep{8},\rep{1})$ is the gluon adjoint representation.
The $(\rep{1},\rep{3})$ is the $W$-boson adjoint representation.

However, here the $\mathrm{SU}(2)$ factor comes from the $k_d = 2$
\emph{massive} directions.  In the SM, $\mathrm{SU}(2)_L$ is the
symmetry of the left-handed doublets, which are massless before
electroweak symmetry breaking.  In the ISS framework, the $k_d$
directions are the ones that have tree-level masses from the
unsatisfied $F$-terms.

This produces a conceptual inversion: the $\mathrm{SU}(2)$ that
looks like $\mathrm{SU}(2)_L$ acts on the
\emph{massive} ISS directions, while in the SM it acts on particles
that are \emph{massless} before EWSB.

\subsection{The tension and Seiberg duality}

Note that $(3,5)$ and $(2,5)$ are \emph{distinct electric theories}
--- $\mathrm{SU}(3)$ SQCD with 5 flavors and $\mathrm{SU}(2)$ SQCD
with 5 flavors respectively.
Standard Seiberg duality maps $\mathrm{SU}(\Nc)$ with $\Nf$ flavors
to $\mathrm{SU}(\Nf-\Nc)$ with $\Nf$ flavors.  Applied here:
the \emph{magnetic dual} of the $(3,5)$ electric theory is an
$\mathrm{SU}(2)$ theory with 5 flavors, and the magnetic dual of
the $(2,5)$ electric theory is an $\mathrm{SU}(3)$ theory with 5
flavors.  The two starting points are thus Seiberg duals of each
other: $(3,5)_{\text{elec}} = (2,5)_{\text{mag}}$ and
$(2,5)_{\text{elec}} = (3,5)_{\text{mag}}$.

The tension between charge assignment and gauge structure is
therefore not a choice between two unrelated theories but a tension
between the electric and magnetic descriptions of the \emph{same}
physics.  The mismatch can be stated precisely.
In the SM:
\begin{itemize}
\item $\mathrm{SU}(2)_L$ acts on massless left-handed doublets.
\item Masses are generated by EWSB, i.e.\ by \emph{breaking}
  $\mathrm{SU}(2)_L$.
\end{itemize}
In the ISS mechanism:
\begin{itemize}
\item $\mathrm{SU}(k_u)$ acts on the \emph{massless} pseudo-moduli.
\item $\mathrm{SU}(k_d)$ acts on the \emph{massive} directions.
\end{itemize}
For $(3,5)$: $\mathrm{SU}(2) = \mathrm{SU}(k_u)$ is the massless
sector, matching the SM role of $\mathrm{SU}(2)_L$.
For $(2,5)$: $\mathrm{SU}(2) = \mathrm{SU}(k_d)$ is the massive
sector, opposite to the SM role of $\mathrm{SU}(2)_L$.

The ``natural'' assignment thus depends on which feature one
prioritises: if the ISS massless/massive split is to match the SM
mass hierarchy, then $(3,5)$ is preferred.  If the non-abelian gauge
group of the residual flavor symmetry is to match the SM gauge group,
then $(2,5)$ is preferred.
The early composite model of Masiero and
Veneziano~\cite{MasieroVeneziano1985}, which studied SQCD with
$N = M = 3$ as a prototype for one lepton family, already faced a
similar tension between composite fermion quantum numbers and the
gauge embedding.


%% ===================================================================
\section{Discussion}
\label{sec:discussion}
%% ===================================================================

\subsection{The sBootstrap choice}

The sBootstrap program~\cite{Rivero2024} adopts $(\Nc,\Nf) = (3,5)$
because the charge assignment takes priority: the $k_u = 2$,
$k_d = 3$ split directly matches the two up-type and three down-type
quarks of the SM (with the top decoupled).
The 25 mesons $M_{ij} = \tilde q_i q_j$ are then identified with
the scalar partners of the SM fermions, and the closure conditions
of the sBootstrap --- the requirement that the composite scalars fill
exactly the sfermion content of one generation --- are satisfied.
The $\mathrm{SU}(2)$ flavor symmetry of the massless sector plays
the role of a proto-electroweak symmetry.

\subsection{The alternative perspective}

The $(2,5)$ starting point has a different appeal: the residual
flavor symmetry is literally
$\mathrm{SU}(3) \times \mathrm{SU}(2) \times \mathrm{U}(1)$, and
the meson representations include $(\rep{3},\rep{2})$,
$(\rep{8},\rep{1})$, and $(\rep{1},\rep{3})$.
If the flavor symmetry could be \emph{gauged} (or if it emerges as a
gauge symmetry in some UV completion), the SM gauge group would
appear automatically.
The price is that the charge assignment is inverted, and the
$\mathrm{SU}(2)$ factor acts on the massive rather than the massless
sector.

\subsection{The $\Nf = 6 \to 5$ transition}

The two perspectives can be connected through the $(3,6)$ starting
point.  With $\Nf = 6$ and $\Nc = 3$, the ISS split is symmetric
($k_u = k_d = 3$), and the theory is self-dual.
Decoupling the heaviest quark (the top) reduces $\Nf$ from 6 to 5.
The resulting $(3,5)$ theory is the sBootstrap case.

Alternatively, one can consider the $(2,6)$ theory and decouple a
flavor to reach $(2,5)$.  However, $(2,6)$ sits at the boundary of
the conformal window, making it a less controlled starting point.

The $\Nf = 6 \to 5$ transition via top decoupling thus provides a
natural connection between the Cacciapaglia--Sannino framework (which
starts at $\Nf = 6$) and the sBootstrap (which works at $\Nf = 5$).
Which of the two $\Nf = 5$ endpoints --- $(3,5)$ or $(2,5)$ ---
is reached depends on the value of $\Nc$ in the electric theory,
i.e.\ on whether the confining dynamics is $\mathrm{SU}(3)$ or
$\mathrm{SU}(2)$.

\subsection{Summary of findings}

The results are:
\begin{enumerate}
\item The ISS rank condition provides a rigid, parameter-free split
  of $\Nf$ flavors into $k_u + k_d$ groups, determined entirely by
  $(\Nc, \Nf)$.
\item Of the four cases examined, only $(\Nc,\Nf) = (3,5)$ gives a
  direct match between the ISS split and the SM charge assignment
  with 2 up-type and 3 down-type flavors.
\item The case $(\Nc,\Nf) = (2,5)$ reproduces the SM gauge group
  $\mathrm{SU}(3) \times \mathrm{SU}(2) \times \mathrm{U}(1)$ from
  the flavor decomposition, but with an inverted charge--mass
  correspondence.
\item The symmetric case $(\Nc,\Nf) = (3,6)$ requires additional
  dynamics (top decoupling) to break the $k_u = k_d$ symmetry and
  reach the SM spectrum.
\item The case $(\Nc,\Nf) = (2,6)$ falls outside the conformal
  window and does not admit a natural SM embedding.
\end{enumerate}

\noindent
The tension between the $(3,5)$ and $(2,5)$ viewpoints is
precisely the tension between the electric and magnetic descriptions
of the same theory under Seiberg duality: $(3,5)$ privileges
the correct charge assignment at the price of a non-SM gauge
structure, while its magnetic dual $(2,5)$ privileges the correct
gauge structure at the price of an inverted mass--charge
correspondence.
The $(3,5)$ theory descends from the $\Nf = 6$ parent by top
decoupling (Section~\ref{sec:36}), and its Seiberg dual is the
$(2,5)$ theory.  The question of which description is
``fundamental'' may ultimately reduce to a question about the
energy scale at which the SM is embedded.

\bigskip
\noindent\textbf{Acknowledgments.}
This work has been partially supported by grant
PID2021-127726NB-I00 funded by MCIN/AEI/10.13039/501100011033.

%% ===================================================================
\begin{thebibliography}{99}

\bibitem{ISS2006}
K.~Intriligator, N.~Seiberg and D.~Shih,
``Dynamical SUSY breaking in meta-stable vacua,''
JHEP \textbf{04} (2006) 021
[hep-th/0602239].

\bibitem{Seiberg1995}
N.~Seiberg,
``Electric-magnetic duality in supersymmetric non-Abelian gauge theories,''
Nucl.\ Phys.\ B \textbf{435} (1995) 129
[hep-th/9411149].

\bibitem{Rivero2024}
A.~Rivero,
``An interpretation of scalars in SO(32),''
Eur.\ Phys.\ J.\ C \textbf{84} (2024) 1058.

\bibitem{CacciapagliaSannino2024}
G.~Cacciapaglia and F.~Sannino,
``Charting Standard Model Duality and its Signatures,''
arXiv:2407.17281 [hep-ph].

\bibitem{MasieroVeneziano1985}
A.~Masiero and G.~Veneziano,
``Split light composite supermultiplets,''
Nucl.\ Phys.\ B \textbf{249} (1985) 593--610.

\end{thebibliography}

\end{document}
